% Part: incompleteness
% Chapter: arithmetization-syntax
% Section: coding-formulas

\documentclass[../../../include/open-logic-section]{subfiles}

\begin{document}

\olfileid{inc}{art}{frm}
\olsection{رمزگذاری \printtoken{P}{formula}}

پس از آنکه رابطهٔ~$\fn{Term}(x)$ را به‌طور بازگشتی اولیه تعریف کردیم،
می‌توانیم از آن برای تعریف بازگشتی اولیهٔ رابطهٔ متناظر برای
!!{formula}s، یعنی~$\fn{Frm}(x)$، استفاده کنیم.

\begin{prop}
رابطهٔ~$\fn{Atom}(x)$ که برقرار است اگر و تنها اگر $x$ عدد گودل یک
!!{formula} اتمی باشد، بازگشتی اولیه است.
\end{prop}

\begin{proof}
عدد~$x$ عدد گودل یک !!{formula} اتمی است اگر و تنها اگر یکی از شرایط
زیر برقرار باشد:
\begin{enumerate}
\item $n,j<x$ و~$z<x$ای وجود داشته باشند به‌گونه‌ای که
  $\len{z}=n$، برای هر~$i<n$، $\fn{Term}((z)_i)$، و $x=$
\[
\Gn{\Obj P^n_j(} \concat \fn{flatten}(z) \concat \Gn{)}.
\]
\item $z_1,z_2<x$ای وجود داشته باشند به‌گونه‌ای که
  $\fn{Term}(z_1)$، $\fn{Term}(z_2)$ و $x=$
\[
\Gn{{\eq}(} \concat z_1 \concat \Gn{,} \concat z_2 \concat{\Gn{)}}.
\]
  \tagitem{prvFalse}{$x = \Gn{\lfalse}$.}{}
  \tagitem{prvTrue}{$x = \Gn{\ltrue}$.}{}
\end{enumerate}
\end{proof}

\begin{prop}
\ollabel{prop:frm-primrec}
رابطهٔ~$\fn{Frm}(x)$ که برقرار است اگر و تنها اگر $x$ عدد گودل یک
!!a{formula} باشد، بازگشتی اولیه است.
\end{prop}

\begin{proof}
دنبالهٔ نمادهای~$s$ یک !!a{formula} است اگر و تنها اگر دنبالهٔ تشکیلی
چون~$s_0$،~\dots، $s_{k-1}=s$ از !!{formula}ها وجود داشته باشد که ثبت
کند $s$ چگونه طبق قواعد تشکیل از !!{formula}sی اتمی ساخته شده است.
می‌توان چنین دنباله‌ای را بدون تکرار برگزید؛ در این صورت هر~$s_i$
زیرفرمولی از~$s$ است، تعداد اعضای دنباله حداکثر~$k$ و عدد گودل
هر عضو حداکثر~$x$ است، که در آن $k=\len{x}$. پس رمز دنبالهٔ تشکیل
$\tuple{s_0,\dots,s_{k-1}}$ را می‌توان با~$p_{k-1}^{k(x+1)}$
کراندار کرد. این یک کران بازگشتی اولیه
برای کمیت‌گذاری بر رمز دنبالهٔ تشکیل به دست می‌دهد.
\end{proof}

\begin{prob}
اثباتی تفصیلی از \olref[inc][art][frm]{prop:frm-primrec} در امتداد
اثبات نخستِ \olref[inc][art][trm]{prop:term-primrec} ارائه کنید.
\end{prob}

\begin{prop}
\ollabel{prop:freeocc-primrec}
رابطهٔ~$\fn{FreeOcc}(x,z,i)$، که برقرار است اگر و تنها اگر نماد
$i$اُمِ فرمول دارای عدد گودل~$x$ رخدادی آزاد از متغیر دارای عدد
گودل~$z$ باشد، بازگشتی اولیه است.
\end{prop}

\begin{proof}
تمرین.
\end{proof}

\begin{prob}
\olref[inc][art][frm]{prop:freeocc-primrec} را ثابت کنید. می‌توانید
از این واقعیت استفاده کنید که هر زیررشتهٔ یک !!a{formula} که خود
!!a{formula} باشد، زیر!!{formula} آن است.
\end{prob}

\begin{prop}
ویژگی~$\fn{Sent}(x)$ که برقرار است اگر و تنها اگر $x$ عدد گودل یک
!!a{sentence} باشد، بازگشتی اولیه است.
\end{prop}

\begin{proof}
!!{sentence} یک !!a{formula} بدون رخداد آزادِ !!{variable}sست. پس
$\fn{Sent}(x)$ برقرار است اگر و تنها اگر
\begin{multline*}
\bforall{i<\len{x}}{\bforall{z<x}{}}\\
(\bexists{j<z}{z=\Gn{\Obj v_j}} \lif \lnot\fn{FreeOcc}(x,z,i)).
\end{multline*}
\end{proof}


\end{document}
